\subsection{Extended Universal Codepoint Escape Access}

The earlier escape forms, such as octal escapes and \texttt{\textbackslash xhh}, are useful for small numeric values. Unicode needs a larger address space. Python therefore supports additional string-literal escape forms for inserting Unicode characters by name or by larger hexadecimal codepoint values.

These are still lexical forms. They are processed while Python reads the source code literal.

\begin{verbatim}
text1 = "π"
text2 = "\N{GREEK SMALL LETTER PI}"
text3 = "\u03C0"
text4 = "\U000003C0"

print(f"text1: {text1!r}")
print(f"text2: {text2!r}")
print(f"text3: {text3!r}")
print(f"text4: {text4!r}")

print()
print(f"text1 == text2: {text1 == text2}")
print(f"text1 == text3: {text1 == text3}")
print(f"text1 == text4: {text1 == text4}")
\end{verbatim}

Possible output:

\begin{verbatim}
text1: 'π'
text2: 'π'
text3: 'π'
text4: 'π'

text1 == text2: True
text1 == text3: True
text1 == text4: True
\end{verbatim}

All four spellings create the same final one-character string.

\subsubsection{Formal Lexical Extraction: Querying Literals via System Naming Maps (\texttt{\textbackslash N\{...\}})}

The escape \texttt{\textbackslash N\{...\}} inserts a Unicode character by its official Unicode name.

\begin{verbatim}
pi = "\N{GREEK SMALL LETTER PI}"
snowman = "\N{SNOWMAN}"
face = "\N{GRINNING FACE}"

print(f"pi:      {pi!r}      U+{ord(pi):04X}")
print(f"snowman: {snowman!r}      U+{ord(snowman):04X}")
print(f"face:    {face!r}      U+{ord(face):04X}")
\end{verbatim}

Possible output:

\begin{verbatim}
pi:      'π'      U+03C0
snowman: '☃'      U+2603
face:    '😀'      U+1F600
\end{verbatim}

The name must be known by Python's Unicode database. The standard library module \texttt{unicodedata} lets us query the same naming system at runtime.

\begin{verbatim}
import unicodedata

chars = ["A", "π", "☃", "😀"]

for ch in chars:
    name = unicodedata.name(ch)
    print(f"{ch!r:<4} U+{ord(ch):06X} {name}")
\end{verbatim}

Possible output:

\begin{verbatim}
'A'  U+000041 LATIN CAPITAL LETTER A
'π'  U+0003C0 GREEK SMALL LETTER PI
'☃'  U+002603 SNOWMAN
'😀' U+01F600 GRINNING FACE
\end{verbatim}

The lookup can also go in the opposite direction.

\begin{verbatim}
import unicodedata

names = [
    "LATIN CAPITAL LETTER A",
    "GREEK SMALL LETTER PI",
    "SNOWMAN",
    "GRINNING FACE",
]

for name in names:
    ch = unicodedata.lookup(name)
    print(f"{name:<30} -> {ch!r} U+{ord(ch):06X}")
\end{verbatim}

Possible output:

\begin{verbatim}
LATIN CAPITAL LETTER A         -> 'A' U+000041
GREEK SMALL LETTER PI          -> 'π' U+0003C0
SNOWMAN                        -> '☃' U+002603
GRINNING FACE                  -> '😀' U+01F600
\end{verbatim}

A bad Unicode name is not a runtime string error. It is a source parsing error. To demonstrate it safely, we compile source text from another string.

\begin{verbatim}
source = 'bad = "\\N{THIS NAME DOES NOT EXIST}"'

try:
    compile(source, "<demo>", "exec")
    print("compiled")
except SyntaxError as err:
    print("syntax error")
    print(f"message: {err.msg}")
\end{verbatim}

Possible output:

\begin{verbatim}
syntax error
message: (unicode error) 'unicodeescape' codec can't decode bytes
in position 0-29: unknown Unicode character name
\end{verbatim}

The exact position numbers may differ if the source string changes. The important rule remains the same: \texttt{\textbackslash N\{...\}} requires a valid Unicode character name.

\subsubsection{Planar Transformations: 16-Bit Base-16 Codepoint Escapes via \texttt{\textbackslash u}}

The escape \texttt{\textbackslash u} inserts one Unicode character by a four-digit hexadecimal codepoint.

The shape is:

\begin{verbatim}
"\uhhhh"
\end{verbatim}

where each \texttt{h} is a hexadecimal digit.

\begin{verbatim}
text = "\u0041\u03C0\u2603"

print(f"text: {text!r}")
print(text)

print()
for ch in text:
    print(f"{ch!r:<4} U+{ord(ch):04X}")
\end{verbatim}

Possible output:

\begin{verbatim}
text: 'Aπ☃'
Aπ☃

'A'  U+0041
'π'  U+03C0
'☃'  U+2603
\end{verbatim}

The escape \texttt{\textbackslash u0041} creates the same character as ordinary \texttt{"A"}.

\begin{verbatim}
normal = "A"
escaped = "\u0041"

print(f"normal:  {normal!r}")
print(f"escaped: {escaped!r}")
print(f"same:    {normal == escaped}")
\end{verbatim}

Possible output:

\begin{verbatim}
normal:  'A'
escaped: 'A'
same:    True
\end{verbatim}

The four hex digits are required. Fewer digits produce a syntax error while the source code is parsed.

\begin{verbatim}
source = 'bad = "\\u03C"'

try:
    compile(source, "<demo>", "exec")
    print("compiled")
except SyntaxError as err:
    print("syntax error")
    print(f"message: {err.msg}")
\end{verbatim}

Possible output:

\begin{verbatim}
syntax error
message: (unicode error) 'unicodeescape' codec can't decode bytes
in position 0-4: truncated \uXXXX escape
\end{verbatim}

The \texttt{\textbackslash u} form reaches codepoints from \texttt{U+0000} through \texttt{U+FFFF}. For characters above that range, use the eight-digit \texttt{\textbackslash U} form.

\begin{verbatim}
pi = "\u03C0"
snowman = "\u2603"

print(f"pi:      {pi!r} U+{ord(pi):04X}")
print(f"snowman: {snowman!r} U+{ord(snowman):04X}")
\end{verbatim}

Possible output:

\begin{verbatim}
pi:      'π' U+03C0
snowman: '☃' U+2603
\end{verbatim}

The practical rule is:

\begin{quote}
\texttt{\textbackslash uhhhh} inserts one Unicode character using exactly four hexadecimal digits.
\end{quote}

\subsubsection{Absolute Space Mapping: 32-Bit Base-16 Extended Codepoint Escapes via \texttt{\textbackslash U}}

The escape \texttt{\textbackslash U} inserts one Unicode character by an eight-digit hexadecimal codepoint.

The shape is:

\begin{verbatim}
"\Uhhhhhhhh"
\end{verbatim}

where each \texttt{h} is a hexadecimal digit.

\begin{verbatim}
face = "\U0001F600"
g_clef = "\U0001D11E"
pi = "\U000003C0"

print(f"face:   {face!r}   U+{ord(face):06X}")
print(f"g_clef: {g_clef!r}   U+{ord(g_clef):06X}")
print(f"pi:     {pi!r}   U+{ord(pi):06X}")
\end{verbatim}

Possible output:

\begin{verbatim}
face:   '😀'   U+01F600
g_clef: '𝄞'   U+01D11E
pi:     'π'   U+0003C0
\end{verbatim}

The eight-digit form can represent codepoints that the four-digit \texttt{\textbackslash u} form cannot write directly.

The number of characters is still counted at the text level.

\begin{verbatim}
text = "\U0001F600\U0001D11E\U000003C0"

print(f"text: {text!r}")
print(f"len(text): {len(text)}")

print()
for ch in text:
    print(f"{ch!r:<4} U+{ord(ch):06X}")
\end{verbatim}

Possible output:

\begin{verbatim}
text: '😀𝄞π'
len(text): 3

'😀' U+01F600
'𝄞' U+01D11E
'π'  U+0003C0
\end{verbatim}

The \texttt{\textbackslash U} escape also requires exactly eight hexadecimal digits.

\begin{verbatim}
source = 'bad = "\\U0001F60"'

try:
    compile(source, "<demo>", "exec")
    print("compiled")
except SyntaxError as err:
    print("syntax error")
    print(f"message: {err.msg}")
\end{verbatim}

Possible output:

\begin{verbatim}
syntax error
message: (unicode error) 'unicodeescape' codec can't decode bytes
in position 0-8: truncated \UXXXXXXXX escape
\end{verbatim}

Unicode also has a maximum valid codepoint: \texttt{U+10FFFF}. A larger value is not accepted.

\begin{verbatim}
source = 'bad = "\\U00110000"'

try:
    compile(source, "<demo>", "exec")
    print("compiled")
except SyntaxError as err:
    print("syntax error")
    print(f"message: {err.msg}")
\end{verbatim}

Possible output:

\begin{verbatim}
syntax error
message: (unicode error) 'unicodeescape' codec can't decode bytes
in position 0-9: illegal Unicode character
\end{verbatim}

The three universal escape forms can be compared directly.

\begin{verbatim}
name_escape = "\N{GREEK SMALL LETTER PI}"
u_escape = "\u03C0"
U_escape = "\U000003C0"

print(f"name_escape: {name_escape!r} U+{ord(name_escape):04X}")
print(f"u_escape:    {u_escape!r} U+{ord(u_escape):04X}")
print(f"U_escape:    {U_escape!r} U+{ord(U_escape):04X}")

print()
print(f"all same: {name_escape == u_escape == U_escape}")
\end{verbatim}

Possible output:

\begin{verbatim}
name_escape: 'π' U+03C0
u_escape:    'π' U+03C0
U_escape:    'π' U+03C0

all same: True
\end{verbatim}

The source spelling differs. The resulting string component is the same.

\begin{quote}
\texttt{\textbackslash N\{...\}} uses a Unicode character name.  
\texttt{\textbackslash uhhhh} uses exactly four hexadecimal digits.  
\texttt{\textbackslash Uhhhhhhhh} uses exactly eight hexadecimal digits and reaches the full valid Unicode codepoint range.
\end{quote}